\documentclass[11pt]{article}
\usepackage{amsmath,amssymb,amsthm,geometry,hyperref}
\geometry{margin=1in}
\newtheorem{theorem}{Theorem}
\newtheorem{lemma}{Lemma}
\newtheorem{corollary}{Corollary}
\setlength{\emergencystretch}{2em}
\title{The osculating mean of the polynomial--logarithmic curve\\
\large A solution of MathDB 327618}
\author{Research notes in Principia Math In Progress}
\date{11 September 2026}
\begin{document}
\maketitle

\section{Task 1: the problem and its solution}
Fix $n\geq2$ and distinct positive numbers $a_1,\ldots,a_n$. For
\[
 \alpha(t)=(t,t^2,\ldots,t^{n-1},\log t),\qquad t>0,
\]
let $O_a$ be the affine hyperplane through $\alpha(a)$ whose direction
contains $\alpha'(a),\ldots,\alpha^{(n-1)}(a)$. If their intersection is
$P=(i_1,\ldots,i_n)$, the last coordinate mean is $M_n=\exp(i_n)$.
Horwitz's final conjecture in \cite{horwitz} asserts that this equals the
multivariate identric mean $I_Z$ defined below. The MathDB title calls this
a logarithmic mean conjecture; the formula to be identified is identric.

Put $N=n-1$, $H_N=\sum_{k=1}^N1/k$, and
\[
 w_i=\frac{a_i^N}{\prod_{k\ne i}(a_i-a_k)}.
\]
\begin{theorem}\label{main}
The hyperplanes $O_{a_1},\ldots,O_{a_n}$ have a unique intersection, and
\[
 i_n=\sum_{i=1}^n w_i\log a_i-H_N,
 \qquad
 M_n=\exp\left(\sum_{i=1}^n w_i\log a_i-H_N\right)=I_Z(a_1,\ldots,a_n).
\]
Thus the source's conjecture holds in every dimension $n\geq2$.
\end{theorem}
The alternating Vandermonde formula in the source is exactly this expression:
if $V(a)=\prod_{j<i}(a_i-a_j)$ and $V_i$ is the same Vandermonde product
after deleting $a_i$, then
\[
 \frac{(-1)^{n+i}V_i(a)}{V(a)}
 =\frac1{\prod_{k\ne i}(a_i-a_k)}.
\]
This identity holds in the displayed input ordering, without a sorting
assumption.

\section{Osculating hyperplanes from Taylor polynomials}
\begin{lemma}\label{plane}
For $a>0$ the hyperplane $O_a$ has equation
\[
 z=\log a-H_N+
 \sum_{j=1}^N\frac{(-1)^{j-1}}{j}\binom Nj a^{-j}x_j.
\]
\end{lemma}
\begin{proof}
The first $N$ coordinates of the vectors
$\alpha'(a),\ldots,\alpha^{(N)}(a)$ form a triangular $N\times N$
matrix with diagonal $1!,2!,\ldots,N!$. They are independent. Hence an
osculating hyperplane exists uniquely and can be expressed as a graph in $z$.

The degree-$N$ Taylor polynomial of $\log t$ about $a$ is
\[
 Q_a(t)=\log a+\sum_{k=1}^N\frac{(-1)^{k-1}}{k a^k}(t-a)^k.
\]
Its constant coefficient is $\log a-H_N$. For $1\leq j\leq N$,
its coefficient of $t^j$ is
\begin{align*}
 \frac{(-1)^{j-1}}{a^j}\sum_{k=j}^N\frac1k\binom kj
 &=\frac{(-1)^{j-1}}{j a^j}\sum_{k=j}^N\binom{k-1}{j-1}\\
 &=\frac{(-1)^{j-1}}{j a^j}\binom Nj.
\end{align*}
The last equality is the elementary telescoping binomial identity.
Replace each $t^j$ in $Q_a$ by the coordinate $x_j$. The resulting graph
has contact through derivative order $N$ with $\alpha$ at $a$, because
$\log t-Q_a(t)$ and its first $N$ derivatives vanish there.
By uniqueness it is $O_a$.
\end{proof}

\section{Interpolation at reciprocal nodes}
\begin{proof}[Proof of Theorem~\ref{main}]
Set $r_i=1/a_i$ and $c_j=(-1)^{j-1}\binom Nj/j$. Define a polynomial
of degree at most $N$ by
\[
 R(r)=z-\sum_{j=1}^N c_j x_j r^j.
\]
All the $c_j$ are nonzero, so specifying its coefficients is equivalent to
specifying $(x_1,\ldots,x_N,z)$. Lemma~\ref{plane} says that membership in
all the $O_{a_i}$ is equivalent to
\[
 R(r_i)=\log a_i-H_N,\qquad 1\leq i\leq n.
\]
The $n=N+1$ nodes $r_i$ are distinct. Ordinary polynomial interpolation
therefore gives exactly one such $R$, proving existence and uniqueness
of the intersection, not merely a formula conditional on its existence.

Its constant coefficient is obtained by evaluating the Lagrange formula
at zero:
\[
 z=R(0)=\sum_{i=1}^n(\log a_i-H_N)
 \prod_{k\ne i}\frac{-r_k}{r_i-r_k}.
\]
Each product is
\[
 \prod_{k\ne i}\frac{-1/a_k}{1/a_i-1/a_k}
 =\prod_{k\ne i}\frac{a_i}{a_i-a_k}=w_i.
\]
The Lagrange basis polynomials sum to one, so $\sum_iw_i=1$.
Consequently $z=\sum_iw_i\log a_i-H_N$, as claimed.
Taking the inverse of the last coordinate function, namely exponentiation,
proves the identity of means.
\end{proof}

\section{Internality and coinciding inputs}
Although the conjecture assumes distinct inputs, the following also checks
that the constructed value is a mean and has the expected continuous extension.
\begin{corollary}\label{internal}
For all positive inputs, including repetitions, the continuous extension is
\[
 M_n(a)=\exp\left(\mathbb E\log\Bigl(\sum_{i=1}^n W_i a_i\Bigr)\right),
\]
where $W$ is uniform on the $(n-1)$-simplex. It is symmetric, continuous,
strictly increasing in each coordinate, and homogeneous of degree one.
It equals $a$ on $(a,\ldots,a)$ and lies strictly between the smallest
and largest input for every nonconstant tuple.
\end{corollary}
\begin{proof}
For $f(t)=t^N\log t$, the sum in Theorem~\ref{main} is the divided
difference $[f](a_1,\ldots,a_n)$. Repeated differentiation gives
\[
 f^{(N)}(t)=N!(\log t+H_N).
\]
The Hermite--Genocchi formula, obtained by repeatedly applying the fundamental
theorem of calculus to divided differences, is
\[
 [f](a_1,\ldots,a_n)=\frac1{N!}
 \mathbb E f^{(N)}\Bigl(\sum_iW_i a_i\Bigr).
\]
This proves the integral formula for distinct inputs. Its right side is
continuous for all positive inputs and therefore gives the unique continuous
extension. Permutation invariance and scaling follow directly. Since
$W_i>0$ almost surely, increasing any one input increases the logarithm
almost surely. For a nonconstant tuple the convex combination lies strictly
between its extreme inputs almost surely, proving strict internality.
\end{proof}
The assertion about repeated inputs concerns the continuous mean. Identical
hyperplanes by themselves do not determine a unique intersection point.

\section{Follow-up tasks 2--5}
\begin{enumerate}
\setcounter{enumi}{1}
\renewcommand{\labelenumi}{\textbf{Task \arabic{enumi}.}}
\item Replace the last coordinate $\log t$ by $t^p$, with real
$p\notin\{0,1,\ldots,n-1\}$. Determine the intersection formula and
classify the parameters for which its last coordinate defines a strict
mean after taking the inverse coordinate function. Recover the logarithmic
case through the limit of $(t^p-1)/p$ as $p\to0$.
\item Construct a geometric analogue for the weighted mean
$\exp(\mathbb E\log\sum_iW_i a_i)$ when $W$ has a prescribed Dirichlet
distribution. Determine which weights can be realized by confluent
osculating conditions and extend the construction beyond integer parameters.
\item Develop an algorithm for the complete intersection point that stays
accurate when the input nodes almost coincide. Prove error bounds and
compare its conditioning and cost with direct Vandermonde inversion.
\item Determine sharp constants bounding
$\log(\bar a/M_n(a))$ by multiples of
$n^{-1}\sum_i(a_i-\bar a)^2/\bar a^2$ under
$\max_i a_i/\min_i a_i\leq R$. Identify extremizers and the asymptotics
as $R\downarrow1$, $R\to\infty$, and $n\to\infty$.
\end{enumerate}

\begin{thebibliography}{9}
\bibitem{horwitz} A. Horwitz, \emph{A logarithmic mean and intersections of
osculating hyperplanes in $\mathbb R^n$}, Journal of Classical Analysis 4
(2014), 41--61, final conjecture on p.~61.
\url{https://doi.org/10.7153/jca-04-02}.
Preprint: \url{https://arxiv.org/abs/1404.1489}.
\bibitem{mathdb} MathDB 327618, live record retrieved 11 September 2026.
\url{https://mathdb.com/p/327618}.
\end{thebibliography}
\end{document}
